\begin{frame}{Bellek Optimizasyonu -- Arduino Uno: 2KB RAM}
  \begin{columns}[T]
    \column{0.48\textwidth}
      \begin{block}{Sorun: String Nesneleri}
        \small
        \texttt{String} sinifi heap'i \textbf{parcalar}:
        \begin{itemize}
          \item Dinamik bellek tahsisi
          \item Heap parcalanmasi
          \item SD baslatma icin yer kalmaz
          \item[\ding{55}] SD kart hatalari
        \end{itemize}
      \end{block}
      \vspace{6pt}
      \begin{block}{Cozum: char{[}{]} + F()}
        \small
        \begin{itemize}
          \item[\ding{51}] \texttt{char uptimeBuf[20]}
          \item[\ding{51}] \texttt{char durumBuf[40]}
          \item[\ding{51}] \texttt{F("sabit metin")}
          \item[] Sabit metinler Flash'a tasindi
          \item[] $\approx$400--500 byte RAM kazanildi
        \end{itemize}
      \end{block}

    \column{0.48\textwidth}
      \begin{block}{Sensor Gürültu Azaltma}
        \small
        \textbf{LM35:} 8 ornek ortalamasi
        \[ T = \frac{1}{8}\sum_{i=1}^{8} x_i \cdot \frac{V_{ref}}{1023} \cdot 100 \]
        Standart sapma $\sigma \to \sigma/\sqrt{8}$\\[6pt]
        \textbf{LDR:} 4 ornek ortalamasi
        \[ L = \frac{1}{4}\sum_{i=1}^{4} x_i \]
      \end{block}
      \vspace{6pt}
      \begin{alertblock}{Kritik Not}
        \small \texttt{F()} makrosu olmadan buyuk
        kodlarda RAM dolarak program kilitleniyor.
      \end{alertblock}
  \end{columns}
\end{frame}
